\documentclass[10pt]{article}
\usepackage[a4paper,margin=0.6in]{geometry}
\usepackage{amsmath,amssymb,amsfonts,mathtools}
\allowdisplaybreaks
\setlength{\jot}{5pt}
\numberwithin{equation}{section}

\title{Cancellation of the three-$\rho$ terms in the finite part of the cross section}
\author{}
\date{}

\begin{document}
\maketitle
\tableofcontents

\section{Overview}

\subsection{The two types of terms}
\begin{center}
\renewcommand{\arraystretch}{1.25}
\begin{tabular}{l|l}
three $\rho$ (explicit factor $i$, or momentum space) & two $\rho$ (no factor $i$)\\ \hline
I.II, \ I.III, \ I.IV, \ I.V first part & I.V second part ($N_{2a}$, $N_{2b\,\rm rest}$)\\
II.I first part ($N_{1a}$, $N_{1b}$) & II.I second part ($2a$--$2f$), \ II.IV\\
III.I, \ III.II, \ III.III first part, \ III.V, \ III.VI & III.III second part
\end{tabular}
\end{center}
``Cancel'' is meant in the following sense. A product of three $\rho$'s is written in one fixed order; two terms cancel if they are equal and opposite as ordered products. The difference between two orderings is a commutator, which contains one $\rho$ less, so what is left after a cancellation is a two-$\rho$ term.

\subsection{What is shown}
\begin{enumerate}
\item \textbf{Cancellation I} (Section~\ref{sec:I}). The six contributions with the kernel $K^{ij}\log\frac{(x-z)^{2}}{(x-w)^{2}}$ --- I.II, I.IV, II.I first part, III.I, III.II, and the first line of your summed Row III --- cancel identically. Only commutators remain, Eq.~\eqref{eq:Omegatot}.
\item \textbf{Cancellation II} (Section~\ref{sec:II}). The $\arccos$ term and the $\log\frac{(z-w)^{2}}{(x-w)^{2}}$ bracket of your summed Row III cancel identically against the part of Row V (first part) in which one of the two projectile gluons does not scatter. Only a commutator and one end-point term remain, Eq.~\eqref{eq:IIresult}.
\item \textbf{What does not cancel} (Section~\ref{sec:left}): (a) the term $-\frac{g^{4}f^{abd}}{8\pi^{4}}K^{ij}$ of Row III that comes from $i\times i\pi$; (b) an end-point term (the last term of $N_{3b}$ and the $\xi\to0$ subtraction of Row V); (c) Row V first part with both gluons scattering; (d) Group III Rows III (first part), V, VI. For (b)--(d) I found no partner in the file with the same phase structure; they are odd under a global reflection for commuting $\rho$ (Section~\ref{sec:refl}), whereas (a) is even.
\item \textbf{The commutators} (Section~\ref{sec:comm}) are two-$\rho$ terms with real coefficients, and the one from Cancellation II has exactly the colour structure of the existing two-$\rho$ term $N_{2a}$ of Row V.
\end{enumerate}

\section{Notation}
\begin{align}
&n^{i}(r)\equiv\frac{r^{i}}{r^{2}},\qquad u^{i}\equiv n^{i}(x'-w'),\qquad v^{j}\equiv n^{j}(y-z),\qquad
\mathcal{P}\equiv\frac{1}{(2\pi)^{3}}\,\frac{ig^{4}}{8\pi^{5}}\,\frac{1}{k^{+}},\notag\\
&a\equiv x-w,\quad b\equiv x-z,\quad c\equiv z-w,\qquad A\equiv a^{2},\ B\equiv b^{2},\ C\equiv c^{2}.
\end{align}
Colour objects:
\begin{align}
\mathcal{D}_{0}^{a}&\equiv\big[U^{ab'}(x')-U^{ab'}(w')\big]\rho^{b'}(x'), &
\Delta U^{bc}(z,y)&\equiv U^{bc}(z)-U^{bc}(y),\notag\\
\mathcal{W}_{\rm III}&\equiv\mathcal{D}_{0}^{a}\,f^{abd}\,U^{de}(x)\,\Delta U^{bc}(z,y)\,\rho^{e}(x)\rho^{c}(y), &
V^{be}&\equiv U^{cb}(z)\,U^{ce}(y),\notag\\
R^{ab}&\equiv f^{cbd}\,U^{ac}(w)\rho^{d}(x)-f^{acd}\,U^{cb}(z)U^{de}(x)\rho^{e}(x), &
\mathcal{Z}&\equiv\mathcal{D}_{0}^{a}\,f^{adc}\,U^{db}(x)\,\Delta U^{ce}(z,y)\,\big[\rho^{e}(y),\rho^{b}(x)\big], \label{eq:colour}
\end{align}
\begin{align}
\mathcal{S}^{a}\equiv f^{hbc}U^{ah}(w)\{\rho^{b}(x),\rho^{c}(y)\}-f^{abc}U^{bd}(x)U^{ce}(y)\{\rho^{d}(x),\rho^{e}(y)\}.
\end{align}
$\mathcal{W}_{\rm III}$ is the colour factor of all terms of Row III. Kernels (as in the Row III note):
\begin{align}
K^{ij}&=\frac{1}{C}\left\{\frac{\delta^{ij}}{2}-\frac{a^{i}a^{j}}{A}-\frac{b^{i}b^{j}}{B}+\frac{2(a\cdot b)\,a^{i}b^{j}}{AB}\right\}, &
\mathcal{A}_{0}&=\int_{0}^{1}\frac{dt}{(a-tb)^{2}}=\frac{\theta_{w}}{\sqrt{AB-(a\cdot b)^{2}}}, &
L_{d}&\equiv\log\frac{C}{A}.
\end{align}

\subsection{Your summed Row III}
The sum written at the end of Row III in the file has five pieces. With $\mathcal{P}$ and the notation above they are
\begin{align}
S_{1}&=\mathcal{P}\int e^{-ik\cdot(w'-w)}\,u^{i}v^{j}\,K^{ij}\log\frac{B}{A}\;\mathcal{W}_{\rm III},\notag\\
S_{2}&=\mathcal{P}\int e^{-ik\cdot(w'-w)}\,\frac{\mathcal{A}_{0}}{C}\Big[(u\cdot c)\,(v\cdot X_{\perp})+(u\cdot Z_{\perp})\,(v\cdot c)\Big]\mathcal{W}_{\rm III},\qquad X_{\perp}=a-\frac{a\cdot b}{B}\,b,\quad Z_{\perp}=b-\frac{a\cdot b}{A}\,a,\notag\\
S_{3}&=\mathcal{P}\int e^{-ik\cdot(w'-w)}\,\frac{L_{d}}{2C}\left[\frac{(u\cdot a)(v\cdot c)}{A}-\frac{(u\cdot c)(v\cdot b)}{B}\right]\mathcal{W}_{\rm III},\notag\\
S_{4}&=-\frac{\mathcal{P}}{2}\int e^{-ik\cdot(w'-w)}\,L_{d}\;u\cdot n(y-w)\;\,n(x-z)\cdot n(w-z)\;\mathcal{D}_{0}^{a}f^{abd}U^{de}(x)\,\Delta U^{bc}(w,y)\,\rho^{e}(x)\rho^{c}(y),\notag\\
S_{5}&=-\frac{1}{(2\pi)^{3}}\,\frac{g^{4}}{8\pi^{4}}\,\frac{1}{k^{+}}\int e^{-ik\cdot(w'-w)}\,u^{i}v^{j}\,K^{ij}\;\mathcal{W}_{\rm III}. \label{eq:S15}
\end{align}
(The prefactor $ig^{4}/16\pi^{5}$ of the file is $\mathcal{P}/2$ without the measure; $\arccos\frac{A-a\cdot b}{\sqrt{AC}}=\theta_{w}$.) In terms of the 22 terms $t_{n}$ of the Row III note, $S_{1}+S_{5}=t_{1}+\dots+t_{15}$, $S_{2}=t_{16}+\dots+t_{19}$, $S_{3}=t_{20}+t_{21}$; I checked numerically that $S_{2}+S_{3}$ equals $\sum_{n=16}^{21}t_{n}$. $S_{5}$ is $\mathcal{P}\times i\pi\,K^{ij}$.

\section{Cancellation I: the $K\log$ terms}\label{sec:I}

\subsection{The six contributions}
All have the kernel $u^{i}v^{j}K^{ij}\log\frac{B}{A}$, the measure $\int_{w,w',x,y,z,x'}$, the phase $e^{-ik\cdot(w'-w)}$, and are of the form $X+{\rm c.c.}$ (for Group II, $N_{1b}=N_{1a}^{*}$). In units of $\mathcal{P}\,u^{i}v^{j}K^{ij}\log\frac{B}{A}$:
\begin{center}
\renewcommand{\arraystretch}{1.3}
\begin{tabular}{p{2.6cm}|l}
contribution & colour factor (order of the $\rho$'s as in the file)\\ \hline
I.II & $\tfrac12\,\mathcal{D}_{0}^{a}\,\mathcal{S}^{a}$\\
I.III ($S_{1}$) & $\mathcal{W}_{\rm III}$\\
I.IV & $\mathcal{D}_{0}^{a}\,V^{be}\rho^{e}(y)\,R^{ab}$\\
II.I first part & $-\tfrac12\big[U^{ac'}(w')\{\rho^{c'}(x'),\rho^{b}(y)\}-U^{ad'}(x')V^{be'}\{\rho^{d'}(x'),\rho^{e'}(y)\}\big]R^{ab}$\\
III.I $+$ III.II & $-\big\{U^{ab'}(x')V^{bd'}\{\rho^{b'}(x'),\rho^{d'}(y)\}-U^{ac'}(w')V^{be'}\rho^{c'}(x')\rho^{e'}(y)-U^{ab'}(x')\rho^{b}(y)\rho^{b'}(x')\big\}R^{ab}$
\end{tabular}
\end{center}
How each row is obtained:
\begin{itemize}
\item \textbf{I.IV}: $U^{bd}(z)U^{be}(y)=V^{de}$; the second term of its bracket is $V^{be}\rho^{e}(y)$ times the second term of $R^{ab}$ because $U^{cb}(z)U^{gb}(z)=\delta^{cg}$.
\item \textbf{III.I, III.II}: $U^{c'd'}(y)U^{c'b}(z)=V^{bd'}$; the two left brackets are added.
\item \textbf{I.II}: the colour factor $\mathcal{S}^{a}$ is antisymmetric under $x\leftrightarrow y$ (relabel $b\leftrightarrow c$, $d\leftrightarrow e$; $f^{hcb}=-f^{hbc}$). On the momenta this exchange is $p\to q-p$, i.e.\ $p^{2}\leftrightarrow l^{2}$ with $l=p-q$. Splitting
\begin{align}
\frac{R^{i}}{p^{2}q^{2}}=\underbrace{\log\frac{q^{2}}{l^{2}}\left[\frac{p^{i}}{p^{2}l^{2}}-\frac{q^{i}}{q^{2}l^{2}}+\frac{q^{i}}{p^{2}q^{2}}\right]}_{\mathcal{R}^{i}}
+2\,\underbrace{\frac{q^{i}}{q^{2}}\left[\frac{1}{l^{2}}\log\frac{q^{2}}{l^{2}}-\frac{1}{p^{2}}\log\frac{p^{2}}{q^{2}}\right]}_{\text{symmetric under }p^{2}\leftrightarrow l^{2}},
\end{align}
the symmetric part drops, and $\mathcal{R}^{i}$ is the Fourier transform of the coordinate kernel,
\begin{align}
\int d^{2}p\,d^{2}q\;e^{iq\cdot(w-y)}e^{ip\cdot(y-x)}\,\mathcal{R}^{i}=-4\pi i\int_{z}v^{j}K^{ij}\log\frac{B}{A},
\qquad\text{so}\qquad
\text{I.II}=\mathcal{P}\int e^{-ik\cdot(w'-w)}u^{i}v^{j}K^{ij}\log\frac{B}{A}\;\tfrac12\mathcal{D}_{0}^{a}\mathcal{S}^{a}.
\end{align}
(Details in the general note. In the new file $R^{i}$ is written with $-2k^{i}$; I read it as $-2q^{i}$.)
\end{itemize}

\subsection{Step 1: terms with $V$}
With $\rho'\equiv\rho(x')$, $\rho_{y}\equiv\rho(y)$ and indices suppressed:
\begin{align}
\text{I.IV}:&\quad\big[U(x')-U(w')\big]\rho'\,V\rho_{y}\,R,\notag\\
\text{II.I}:&\quad+\tfrac12\,U(x')\,V\{\rho',\rho_{y}\}\,R,\notag\\
\text{III}:&\quad-U(x')\,V\{\rho',\rho_{y}\}\,R+U(w')\,V\rho'\rho_{y}\,R.
\end{align}
\begin{itemize}
\item $U(w')V$: $\;-U(w')V\rho'\rho_{y}R+U(w')V\rho'\rho_{y}R=0$.
\item $U(x')V$: $\;U(x')V\big[\rho'\rho_{y}-\tfrac12\{\rho',\rho_{y}\}\big]R=\tfrac12\,U(x')V[\rho',\rho_{y}]R$.
\end{itemize}

\subsection{Step 2: terms without $V$ from II and III}
\begin{align}
-\tfrac12\,U(w')\{\rho',\rho_{y}\}R+U(x')\rho_{y}\rho'R
&=-U(w')\rho'\rho_{y}R+\tfrac12U(w')[\rho',\rho_{y}]R+U(x')\rho'\rho_{y}R-U(x')[\rho',\rho_{y}]R\notag\\
&=\mathcal{D}_{0}^{a}\,\rho^{b}(y)\,R^{ab}-\Big(U^{ab'}(x')-\tfrac12U^{ab'}(w')\Big)\big[\rho^{b'}(x'),\rho^{b}(y)\big]R^{ab}.
\end{align}

\subsection{Step 3: add Row III}
Write $R^{ab}=R_{1}^{ab}+R_{2}^{ab}$ with $R_{1}^{ab}=f^{cbd}U^{ac}(w)\rho^{d}(x)$, $R_{2}^{ab}=-f^{acd}U^{cb}(z)U^{de}(x)\rho^{e}(x)$. Relabelling $c\leftrightarrow b$,
\begin{align}
R_{2}^{ab}\rho^{b}(y)=-f^{abd}\,U^{bc}(z)\,U^{de}(x)\,\rho^{e}(x)\rho^{c}(y),
\end{align}
which is minus the $U(z)$ term of $\mathcal{W}_{\rm III}$. Therefore
\begin{align}
\mathcal{D}_{0}^{a}\rho^{b}(y)R^{ab}+\mathcal{W}_{\rm III}
&=\mathcal{D}_{0}^{a}\Big\{\rho^{b}(y)R_{1}^{ab}+\underbrace{\rho^{b}(y)R_{2}^{ab}-R_{2}^{ab}\rho^{b}(y)}_{[\rho^{b}(y),R_{2}^{ab}]}-f^{abd}U^{de}(x)U^{bc}(y)\rho^{e}(x)\rho^{c}(y)\Big\}.
\end{align}
The $U(z)$ dependence has cancelled up to a commutator. The five lower rows of the table add to $\Omega_{0}+\Omega_{\rm comm}$ with
\begin{align}
\Omega_{0}&=\mathcal{D}_{0}^{a}\Big\{f^{cbd}\,U^{ac}(w)\,\rho^{b}(y)\rho^{d}(x)-f^{abd}\,U^{bc}(y)\,U^{de}(x)\,\rho^{e}(x)\rho^{c}(y)\Big\},\\
\Omega_{\rm comm}&=\tfrac12\,U^{ab'}(x')\,V^{bd'}\big[\rho^{b'}(x'),\rho^{d'}(y)\big]R^{ab}
-\Big(U^{ab'}(x')-\tfrac12U^{ab'}(w')\Big)\big[\rho^{b'}(x'),\rho^{b}(y)\big]R^{ab}\notag\\
&\quad-\mathcal{D}_{0}^{a}\,f^{acd}\,U^{cb}(z)\,U^{de}(x)\,\big[\rho^{b}(y),\rho^{e}(x)\big].
\end{align}

\subsection{Step 4: add Row II}
Relabel the indices of $\Omega_{0}$: in the first term $(c,b,d)\to(h,c,b)$, in the second $(b,d,c,e)\to(c,b,e,d)$:
\begin{align}
\Omega_{0}=\mathcal{D}_{0}^{a}\Big\{-f^{hbc}\,U^{ah}(w)\,\rho^{c}(y)\rho^{b}(x)+f^{abc}\,U^{bd}(x)\,U^{ce}(y)\,\rho^{d}(x)\rho^{e}(y)\Big\}.
\end{align}
Adding $\tfrac12\mathcal{D}_{0}^{a}\mathcal{S}^{a}=\tfrac12\mathcal{D}_{0}^{a}\{f^{hbc}U^{ah}(w)(\rho^{b}(x)\rho^{c}(y)+\rho^{c}(y)\rho^{b}(x))-f^{abc}U^{bd}(x)U^{ce}(y)(\rho^{d}(x)\rho^{e}(y)+\rho^{e}(y)\rho^{d}(x))\}$:
\begin{align}
\Omega_{0}+\tfrac12\,\mathcal{D}_{0}^{a}\mathcal{S}^{a}=\tfrac12\,\mathcal{D}_{0}^{a}\Big\{f^{hbc}\,U^{ah}(w)\big[\rho^{b}(x),\rho^{c}(y)\big]+f^{abc}\,U^{bd}(x)\,U^{ce}(y)\big[\rho^{d}(x),\rho^{e}(y)\big]\Big\}.
\end{align}

\subsection{Result of Cancellation I}
\begin{align}
\boxed{\;\begin{aligned}
&\text{I.II}+S_{1}+\text{I.IV}+\text{II.I(first)}+\text{III.I}+\text{III.II}=\mathcal{P}\int_{w,w',x,y,z,x'}e^{-ik\cdot(w'-w)}\,u^{i}v^{j}K^{ij}\log\frac{(x-z)^{2}}{(x-w)^{2}}\;\Omega_{\rm tot}+{\rm c.c.},\\
&\Omega_{\rm tot}=\Omega_{\rm comm}+\tfrac12\,\mathcal{D}_{0}^{a}\Big\{f^{hbc}\,U^{ah}(w)\big[\rho^{b}(x),\rho^{c}(y)\big]+f^{abc}\,U^{bd}(x)\,U^{ce}(y)\big[\rho^{d}(x),\rho^{e}(y)\big]\Big\}.
\end{aligned}\;} \label{eq:Omegatot}
\end{align}
Every three-$\rho$ term has cancelled; $\Omega_{\rm tot}$ consists of commutators only.

\section{Cancellation II: Row III ($S_{2}+S_{3}$) against Row V, first part}\label{sec:II}

\subsection{Step 1: Row V first part, third term equals the second}
Row V (first part) is, with $\xi=p^{+}/k^{+}$ (resp.\ $q^{+}/k^{+}$) and $\frac{ig^{4}}{16\pi^{4}}\frac{1}{2\pi k^{+}}=\frac{\mathcal{P}}{4}\times(2\pi)^{3}$,
\begin{align}
N_{\rm V(1)}&=\frac{\mathcal{P}}{4}\int_{x,z,y,x',y',w'}\int_{0}^{1}d\xi\;e^{-ik\cdot w'}\Big\{e^{ik\cdot[(1-\xi)z+\xi x]}\Big(\mathcal{K}_{1}-\left[\tfrac1\xi\right]_{+}\mathcal{K}_{2}\Big)\notag\\
&\hspace{4.6cm}-\left[\tfrac1\xi\right]_{+}e^{ik\cdot[(1-\xi)x+\xi z]}\,\mathcal{K}_{3}\Big\}\;\mathcal{D}_{0}^{a}\,\mathcal{C}_{V}^{a}+{\rm c.c.},\notag\\
\mathcal{K}_{1}&=u\cdot n(z-x)\;\,n(y-z)\cdot n(y'-x),\qquad
\mathcal{K}_{2}=n(y'-x)\cdot n(z-x)\;\,u\cdot n(y-z),\notag\\
\mathcal{K}_{3}&=u\cdot n(y'-x)\;\,n(z-x)\cdot n(y-z),\notag\\
\mathcal{C}_{V}^{a}&=f^{adc}\,U^{db}(x)\,U^{ce}(z)\{\rho^{e}(y),\rho^{b}(y')\}-f^{abc}\,U^{cd}(y)\,U^{be}(y')\{\rho^{d}(y),\rho^{e}(y')\}.
\end{align}
Under the exchange $(x,y')\leftrightarrow(z,y)$:
\begin{itemize}
\item $\mathcal{C}_{V}^{a}\to-\mathcal{C}_{V}^{a}$ (relabel $d\leftrightarrow c$, $b\leftrightarrow e$ in the first term and $b\leftrightarrow c$, $d\leftrightarrow e$ in the second);
\item $e^{ik\cdot[(1-\xi)z+\xi x]}\mathcal{K}_{2}\to e^{ik\cdot[(1-\xi)x+\xi z]}\,n(y-z)\cdot n(x-z)\;u\cdot n(y'-x)=-e^{ik\cdot[(1-\xi)x+\xi z]}\mathcal{K}_{3}$.
\end{itemize}
So the third term is the image of the second under a relabelling of integration variables, and
\begin{align}
N_{\rm V(1)}=\frac{\mathcal{P}}{4}\int\int_{0}^{1}d\xi\;e^{-ik\cdot w'}e^{ik\cdot[(1-\xi)z+\xi x]}\Big(\mathcal{K}_{1}-2\left[\tfrac1\xi\right]_{+}\mathcal{K}_{2}\Big)\,\mathcal{D}_{0}^{a}\,\mathcal{C}_{V}^{a}+{\rm c.c.} \label{eq:V1}
\end{align}

\subsection{Step 2: split the colour factor}
Insert $U(x)=\Delta U(x,y')+U(y')$ and $U(z)=\Delta U(z,y)+U(y)$ in the first term of $\mathcal{C}_{V}^{a}$. The piece
\begin{align}
f^{adc}\,U^{db}(y')\,U^{ce}(y)\{\rho^{e}(y),\rho^{b}(y')\}
\end{align}
is identical to the second term of $\mathcal{C}_{V}^{a}$ (relabel $b\to d$, $d\to e$, $e\to b$ there) and cancels it. Exactly:
\begin{align}
\mathcal{C}_{V}^{a}=f^{adc}\Big\{\underbrace{U^{db}(y')\,\Delta U^{ce}(z,y)}_{\rm (i)}+\underbrace{\Delta U^{db}(x,y')\,U^{ce}(y)}_{\rm (ii)}+\underbrace{\Delta U^{db}(x,y')\,\Delta U^{ce}(z,y)}_{\rm (iii)}\Big\}\{\rho^{e}(y),\rho^{b}(y')\}. \label{eq:CVsplit}
\end{align}
In (i) only the gluon from $y$ scatters, in (ii) only the gluon from $y'$, in (iii) both. Pieces (i) and (ii) are the ones that cancel against Row III.

\subsection{Step 3: rename the variables of (i) and (ii)}
Row III has the colour factor $\mathcal{W}_{\rm III}$: a source at $x$ with $U(x)$ and a gluon at $z$ emitted from $y$.
\paragraph{Piece (i).} Rename $y'\to x$ and call the old $x$ (which now appears only in the kernel) $r$:
\begin{align}
\mathcal{D}_{0}^{a}f^{adc}U^{db}(x)\Delta U^{ce}(z,y)\{\rho^{e}(y),\rho^{b}(x)\}
&=2\,\mathcal{D}_{0}^{a}f^{adc}U^{db}(x)\Delta U^{ce}(z,y)\rho^{b}(x)\rho^{e}(y)+\mathcal{Z}=-2\,\mathcal{W}_{\rm III}+\mathcal{Z},
\end{align}
using $\mathcal{W}_{\rm III}=\mathcal{D}_{0}^{a}f^{acd}U^{db}(x)\Delta U^{ce}(z,y)\rho^{b}(x)\rho^{e}(y)$ (relabel $b\to c$, $e\to b$, $c\to e$ in \eqref{eq:colour}) and $f^{acd}=-f^{adc}$. The kernels become
$\mathcal{K}_{1}\to u\cdot n(z-r)\;v\cdot n(x-r)$, $\ \mathcal{K}_{2}\to n(x-r)\cdot n(z-r)\;u\cdot v$, with phase $e^{ik\cdot[(1-\xi)z+\xi r]}$.
\paragraph{Piece (ii).} Rename $y\to x$, $x\to z$, $y'\to y$, and call the old $z$ now $r$. Relabelling $d\leftrightarrow c$, $b\leftrightarrow e$:
\begin{align}
\mathcal{D}_{0}^{a}f^{adc}\Delta U^{db}(z,y)U^{ce}(x)\{\rho^{e}(x),\rho^{b}(y)\}=-\mathcal{D}_{0}^{a}f^{adc}U^{db}(x)\Delta U^{ce}(z,y)\{\rho^{b}(x),\rho^{e}(y)\}=+2\,\mathcal{W}_{\rm III}-\mathcal{Z}.
\end{align}
The kernels become $\mathcal{K}_{1}\to u\cdot n(r-z)\;n(x-r)\cdot v=-u\cdot n(z-r)\;v\cdot n(x-r)$, $\ \mathcal{K}_{2}\to v\cdot n(r-z)\;u\cdot n(x-r)=-v\cdot n(z-r)\;u\cdot n(x-r)$, with phase $e^{ik\cdot[(1-\xi)r+\xi z]}$.
\paragraph{Sum of (i) and (ii).} Substituting $\xi\to1-\xi$ in the $\mathcal{K}_{1}$ term of (ii),
\begin{align}
N_{\rm V(1)}^{\rm (i)+(ii)}&=\mathcal{P}\int e^{-ik\cdot w'}\Big(\mathcal{W}_{\rm III}-\tfrac12\mathcal{Z}\Big)\Big\{-\mathcal{I}_{1}+\mathcal{I}_{2}+\mathcal{I}_{3}\Big\}+{\rm c.c.},\notag\\
\mathcal{I}_{1}&=\int_{r}\int_{0}^{1}d\xi\;e^{ik\cdot[(1-\xi)z+\xi r]}\;u\cdot n(z-r)\;\,v\cdot n(x-r),\notag\\
\mathcal{I}_{2}&=\int_{r}\int_{0}^{1}d\xi\left[\tfrac1\xi\right]_{+}e^{ik\cdot[(1-\xi)r+\xi z]}\;v\cdot n(z-r)\;\,u\cdot n(x-r),\notag\\
\mathcal{I}_{3}&=\int_{r}\int_{0}^{1}d\xi\left[\tfrac1\xi\right]_{+}e^{ik\cdot[(1-\xi)z+\xi r]}\;n(x-r)\cdot n(z-r)\;\,u\cdot v. \label{eq:I123}
\end{align}
(The factor: $\frac{\mathcal{P}}{4}\,(-2\mathcal{W}_{\rm III}+\mathcal{Z})\times2\mathcal{K}_{1}=-\mathcal{P}(\mathcal{W}_{\rm III}-\frac12\mathcal{Z})\mathcal{K}_{1}$, and $\frac{\mathcal{P}}{4}(-2\mathcal{W}_{\rm III}+\mathcal{Z})(-2)[\frac1\xi]_{+}\mathcal{K}_{2}=+\mathcal{P}(\mathcal{W}_{\rm III}-\frac12\mathcal{Z})[\frac1\xi]_{+}\mathcal{K}_{2}$ for (i), and the same for (ii).)

\subsection{Step 4: change of variable $r\to w$}
In Row III the phase is $e^{ik\cdot w}$. In $\mathcal{I}_{1}$ put $w\equiv(1-\xi)z+\xi r$, i.e.\ $r=z+(w-z)/\xi$, $d^{2}r=d^{2}w/\xi^{2}$. With $t\equiv1-\xi$:
\begin{align}
z-r=\frac{c}{\xi}\;\Rightarrow\;n(z-r)=\xi\,\frac{c}{C},\qquad
x-r=b+\frac{c}{\xi}=\frac{a-t\,b}{\xi}\;\Rightarrow\;n(x-r)=\xi\,\frac{a-tb}{(a-tb)^{2}} .
\end{align}
The factors of $\xi$ cancel the Jacobian:
\begin{align}
\mathcal{I}_{1}=\int_{w}e^{ik\cdot w}\int_{0}^{1}dt\;\frac{(u\cdot c)\;v\cdot(a-tb)}{C\,(a-tb)^{2}} .
\end{align}
In $\mathcal{I}_{2}$ put $w\equiv(1-\xi)r+\xi z$; the same formulas hold with $\xi\to1-\xi$, so $n(z-r)=(1-\xi)c/C$, $n(x-r)=(1-\xi)(a-\xi b)/(a-\xi b)^{2}$, and at $\xi=0$ one has $w=r$, so the plus prescription acts on the $\xi$ dependence at fixed $w$:
\begin{align}
\mathcal{I}_{2}=\int_{w}e^{ik\cdot w}\,\frac{v\cdot c}{C}\int_{0}^{1}\frac{d\xi}{\xi}\left[\frac{u\cdot(a-\xi b)}{(a-\xi b)^{2}}-\frac{u\cdot a}{A}\right].
\end{align}

\subsection{Step 5: the one-parameter integrals}
With $(a-tb)^{2}=A-2t\,a\cdot b+t^{2}B$:
\begin{align}
\int_{0}^{1}\frac{dt}{(a-tb)^{2}}=\mathcal{A}_{0},\qquad
\int_{0}^{1}\frac{t\,dt}{(a-tb)^{2}}=\frac{1}{2B}\int_{0}^{1}dt\,\frac{d}{dt}\log(a-tb)^{2}+\frac{a\cdot b}{B}\,\mathcal{A}_{0}=\frac{L_{d}}{2B}+\frac{a\cdot b}{B}\,\mathcal{A}_{0}.
\end{align}
\paragraph{$\mathcal{I}_{1}$.}
\begin{align}
\int_{0}^{1}dt\,\frac{(u\cdot c)\,v\cdot(a-tb)}{C\,(a-tb)^{2}}
&=\frac{u\cdot c}{C}\left[(v\cdot a)\,\mathcal{A}_{0}-(v\cdot b)\left(\frac{L_{d}}{2B}+\frac{a\cdot b}{B}\mathcal{A}_{0}\right)\right]\notag\\
&=u^{i}v^{j}\left[\mathcal{A}_{0}\frac{c^{i}a^{j}}{C}-\mathcal{A}_{0}\frac{(a\cdot b)\,c^{i}b^{j}}{CB}-\frac12L_{d}\frac{c^{i}b^{j}}{CB}\right]=u^{i}v^{j}\,\big(t_{16}+t_{18}+t_{20}\big)^{ij}.
\end{align}
\paragraph{$\mathcal{I}_{2}$.}
\begin{align}
\frac1\xi\left[\frac{a^{i}-\xi b^{i}}{(a-\xi b)^{2}}-\frac{a^{i}}{A}\right]
=\frac{-A\,b^{i}+a^{i}\,(2\,a\cdot b-\xi B)}{A\,(a-\xi b)^{2}},
\end{align}
\begin{align}
\int_{0}^{1}\frac{d\xi}{\xi}\left[\frac{a^{i}-\xi b^{i}}{(a-\xi b)^{2}}-\frac{a^{i}}{A}\right]
&=\left(-b^{i}+\frac{2(a\cdot b)}{A}a^{i}\right)\mathcal{A}_{0}-\frac{B}{A}\,a^{i}\left(\frac{L_{d}}{2B}+\frac{a\cdot b}{B}\mathcal{A}_{0}\right)
=-\mathcal{A}_{0}\,b^{i}+\mathcal{A}_{0}\frac{(a\cdot b)}{A}a^{i}-\frac{L_{d}}{2A}\,a^{i},
\end{align}
so that
\begin{align}
\frac{v\cdot c}{C}\int_{0}^{1}\frac{d\xi}{\xi}[\dots]=-u^{i}v^{j}\left[\mathcal{A}_{0}\frac{b^{i}c^{j}}{C}-\mathcal{A}_{0}\frac{(a\cdot b)\,a^{i}c^{j}}{CA}+\frac12L_{d}\frac{a^{i}c^{j}}{CA}\right]=-u^{i}v^{j}\,\big(t_{17}+t_{19}+t_{21}\big)^{ij}.
\end{align}
\paragraph{Together.}
\begin{align}
-\mathcal{I}_{1}+\mathcal{I}_{2}&=-\int_{w}e^{ik\cdot w}\,u^{i}v^{j}\sum_{n=16}^{21}t_{n}^{ij}\notag\\
&=-\int_{w}e^{ik\cdot w}\left\{\frac{\mathcal{A}_{0}}{C}\Big[(u\cdot c)(v\cdot X_{\perp})+(u\cdot Z_{\perp})(v\cdot c)\Big]+\frac{L_{d}}{2C}\left[\frac{(u\cdot a)(v\cdot c)}{A}-\frac{(u\cdot c)(v\cdot b)}{B}\right]\right\},
\end{align}
which is exactly minus the kernel of $S_{2}+S_{3}$ in \eqref{eq:S15}.

\subsection{Result of Cancellation II}
\begin{align}
\boxed{\;S_{2}+S_{3}+N_{\rm V(1)}^{\rm (i)+(ii)}
=\mathcal{P}\int e^{-ik\cdot(w'-w)}\,u^{i}v^{j}\sum_{n=16}^{21}t_{n}^{ij}\;\tfrac12\,\mathcal{Z}
\;+\;\mathcal{P}\int e^{-ik\cdot w'}\,\mathcal{I}_{3}\,\Big(\mathcal{W}_{\rm III}-\tfrac12\mathcal{Z}\Big)+{\rm c.c.}\;} \label{eq:IIresult}
\end{align}
The three-$\rho$ parts of the $\arccos$ term and of the $\log\frac{(z-w)^{2}}{(x-w)^{2}}$ bracket have cancelled. What is left is a commutator term and the end-point term $\mathcal{I}_{3}$.

\section{What is left of the three-$\rho$ terms}\label{sec:left}

\subsection{(a) The $\pi$ term $S_{5}$}
\begin{align}
S_{5}=-\frac{1}{(2\pi)^{3}}\,\frac{g^{4}}{8\pi^{4}}\,\frac{1}{k^{+}}\int e^{-ik\cdot(w'-w)}\,u^{i}v^{j}K^{ij}\,\mathcal{W}_{\rm III}+{\rm c.c.}
\end{align}
It has the kernel $K^{ij}$ of Cancellation I but only the colour factor of Row III. In Cancellation I the colour factors of the six contributions add to zero (up to commutators). Therefore: \emph{if the logarithm of all six contributions were $\log\frac{-(x-z)^{2}}{(x-w)^{2}}=\log\frac{B}{A}+i\pi$, the $\pi$ terms would cancel in exactly the same way as the logarithms.} With the file as it is (only Row III has the minus sign), $S_{5}$ has nothing to cancel against. This is the place to check in the derivation of I.II, I.IV, II.I, III.I, III.II.

\subsection{(b) The end-point term}
Doing the $\xi$ integral in $\mathcal{I}_{3}$, with $\int_{0}^{1}\frac{d\xi}{\xi}(e^{i\xi\sigma}-1)=-{\rm Ein}(-i\sigma)$,
\begin{align}
\mathcal{I}_{3}=-\,e^{ik\cdot z}\;u\cdot v\int_{r}n(x-r)\cdot n(z-r)\;{\rm Ein}\big(-ik\cdot(r-z)\big).
\end{align}
Renaming $w\to z$, $z\to r$ in $S_{4}$ brings it to the same form (its colour factor becomes $\mathcal{W}_{\rm III}$), and
\begin{align}
S_{4}+\mathcal{P}\int e^{-ik\cdot w'}\mathcal{I}_{3}\,\mathcal{W}_{\rm III}
=-\mathcal{P}\int e^{-ik\cdot(w'-z)}\,u\cdot v\;\mathcal{W}_{\rm III}\int_{r}n(x-r)\cdot n(z-r)\left\{{\rm Ein}\big(-ik\cdot(r-z)\big)+\frac12\log\frac{(r-z)^{2}}{(x-z)^{2}}\right\}.
\end{align}
The two do not cancel. Moreover $\mathcal{W}_{\rm III}$ does not depend on $r$, $n(x-r)\cdot n(z-r)\to1/r^{2}$ and the bracket grows like $\log r^{2}$, so the $r$ integral does not converge at large $r$. In $\mathcal{I}_{3}$ this comes from the subtraction at $\xi=0$ of the plus prescription, whose coefficient $\int_{r}n(x-r)\cdot n(z-r)$ is itself divergent: the $\xi\to0$ limit and the $r$ integral do not commute. (With the last term of $N_{3b}$ as in the earlier version of the file, $-\tfrac12L_{d}\,\delta^{ij}/C$, the finite part of $\mathcal{I}_{3}$ in the variable $w$ is $+\tfrac12L_{d}\,\delta^{ij}/C$ and cancels it, leaving only the divergent piece $\int_{0}^{1}\frac{d\xi}{\xi}\big[\int_{w}\frac{e^{ik\cdot w}}{(z-w)^{2}}-e^{ik\cdot z}\int_{r}n(x-r)\cdot n(z-r)\big]$.)

Note also that the last term of $N_{3b}$ is written with $\Delta U^{bc}(z,y)$ in $N_{3b}$ but with $\Delta U^{bc}(w,y)$ in the summed formula; I used the latter, which is the one consistent with the factor $n(y-w)$.

\subsection{(c) Row V first part, both gluons scattering}
Piece (iii) of \eqref{eq:CVsplit} in \eqref{eq:V1}:
\begin{align}
N_{\rm V(1)}^{\rm (iii)}=\frac{\mathcal{P}}{4}\int\int_{0}^{1}d\xi\;e^{-ik\cdot w'}e^{ik\cdot[(1-\xi)z+\xi x]}\Big(\mathcal{K}_{1}-2\left[\tfrac1\xi\right]_{+}\mathcal{K}_{2}\Big)\,\mathcal{D}_{0}^{a}f^{adc}\Delta U^{db}(x,y')\Delta U^{ce}(z,y)\{\rho^{e}(y),\rho^{b}(y')\}+{\rm c.c.}
\end{align}

\subsection{(d) Group III, Rows III (first part) $+$ V $+$ VI}
These three add to one term with the colour factor (general note)
\begin{align}
\mathcal{C}_{3}(q)&=2f^{c'd'a}\,\mathcal{D}^{c'}_{q}\,\Delta U^{d'b}(z,x)\rho^{b}(x)\,\Delta U^{ac}(w,y)\rho^{c}(y)\notag\\
&\quad+f^{c'd'a}\mathcal{D}^{c'}_{q}\Big\{U^{ac}(w)U^{d'b}(z)-2U^{ac}(y)U^{d'b}(z)+U^{ac}(y)U^{d'b}(x)\Big\}\big[\rho^{c}(y),\rho^{b}(x)\big],
\end{align}
$\mathcal{D}^{c'}_{q}=\Delta U^{c'e'}(q,x')\rho^{e'}(x')$. After renaming, its three-$\rho$ part has the same colour structure as (c) with the opposite sign ($-2f^{adc}\mathcal{D}_{0}^{a}\Delta U\rho\,\Delta U\rho$): (c) and (d) are the virtual and real versions of the same vertex. Their phases differ, $e^{-ik\cdot w'}e^{ik\cdot[(1-\xi)z+\xi x]}$ with $0<\xi<1$ against $e^{-ik\cdot(1+\eta)w'}e^{ik\cdot z}e^{i\eta k\cdot x}$ with $0<\eta<\infty$, and $\mathcal{D}_{0}$ depends on $w'$, so they do not cancel point by point.

\subsection{Reflection property of (b), (c), (d)}\label{sec:refl}
Each of (b), (c), (d) is of the form $i\,G\,\mathcal{W}+{\rm c.c.}$ where the kernel $G$ (phases included) satisfies $G(-{\rm all\ coordinates})=G^{*}$: the phases are conjugated, the products of an even number of $n$'s are even, and ${\rm Ein}(i\sigma)$, $J_{0}(s)$, $\log(is)$, $1/(s-i0)$ go into their conjugates. For commuting $\rho$ the colour factor $\mathcal{W}$ is real and
\begin{align}
i\,G\,\mathcal{W}+{\rm c.c.}=-2\,\mathcal{W}\;{\rm Im}\,G,\qquad{\rm Im}\,G\ \text{odd under the reflection}.
\end{align}
So these terms vanish after the transverse integrations whenever the averages over projectile and target are reflection symmetric. This is not a point-by-point cancellation and it is an assumption about the averages. The same argument applied to $S_{5}$ gives the opposite conclusion: its kernel $i\cdot i\pi K^{ij}e^{-ik\cdot(w'-w)}$ has an \emph{even} real part $-\pi K^{ij}\cos[k\cdot(w'-w)]$ and does not vanish. $S_{5}$ is therefore the one three-$\rho$ term that neither cancels nor vanishes by symmetry.

\section{The two-$\rho$ terms produced by the cancellations}\label{sec:comm}

The commutators left by Cancellations I and II are
\begin{align}
\mathcal{P}\int e^{-ik\cdot(w'-w)}\,u^{i}v^{j}\left\{K^{ij}\log\frac{B}{A}\;\Omega_{\rm tot}+\frac12\sum_{n=16}^{21}t_{n}^{ij}\;\mathcal{Z}\right\}+{\rm c.c.},\qquad
\sum_{n=16}^{21}t_{n}^{ij}=\mathcal{A}_{0}T^{ij}-\frac12L_{d}\left(S^{ij}-\frac{\delta^{ij}}{C}\right).
\end{align}
If $[\rho^{a}(x),\rho^{b}(y)]=i\kappa f^{abc}\rho^{c}(x)\delta^{2}(x-y)$ ($\kappa$ depends on the normalisation of $\rho$), each commutator gives a factor $i$ which with the $i$ of $\mathcal{P}$ makes the coefficient real, and removes one $\rho$ and one integration. For $\mathcal{Z}$, using $f^{adc}U^{db}U^{ce}=U^{ah}f^{hbe}$ and $f^{hbe}f^{ebg}=-N_{c}\delta^{hg}$:
\begin{align}
\mathcal{Z}=i\kappa\;\mathcal{D}_{0}^{a}\Big[-f^{adc}f^{gbe}\,U^{db}(x)\,U^{ce}(z)+N_{c}\,U^{ag}(x)\Big]\rho^{g}(x)\;\delta^{2}(x-y).
\end{align}
This is the colour structure $\big[U(x')-U(w')\big]\big[-f^{adc}f^{fbe}U^{db}(x)U^{ce}(z)+N_{c}U^{af}(x)\big]\rho(x')\rho$ of the two-$\rho$ term $N_{2a}$ of Row V (second part). Similarly the last bracket of $\Omega_{\rm tot}$ becomes $i\kappa N_{c}\,[U^{ae}(w)+U^{ae}(x)]\rho^{e}(x)\delta^{2}(x-y)$.

\section{Summary and points to check}

\begin{center}
\renewcommand{\arraystretch}{1.3}
\begin{tabular}{p{7.2cm}|p{9.2cm}}
three-$\rho$ terms & status\\ \hline
I.II, I.IV, II.I first part, III.I, III.II, Row III $S_{1}$ & cancel identically; commutators $\Omega_{\rm tot}$ remain\\
Row III $S_{2}$ ($\arccos$), $S_{3}$ ($\log$ bracket); Row V first part, pieces (i), (ii) & cancel identically; commutator $\tfrac12\mathcal{Z}$ and end-point term remain\\
Row III $S_{5}$ ($-\frac{g^{4}}{8\pi^{4}}K$) & not cancelled; even under reflection\\
Row III $S_{4}$ and end-point term $\mathcal{I}_{3}$ of Row V & not cancelled; $r$ integral divergent at large $r$\\
Row V first part, piece (iii); III.III first part, III.V, III.VI & not cancelled point by point; odd under reflection for commuting $\rho$
\end{tabular}
\end{center}
Points to check in the derivation:
\begin{enumerate}
\item whether the logarithm in I.II, I.IV, II.I, III.I, III.II should also be $\log\frac{-(x-z)^{2}}{(x-w)^{2}}$; if so $S_{5}$ cancels with them;
\item the last term of $N_{3b}$ (colour factor $U(z)$ or $U(w)$; sign and form), and the order of the $\Lambda\to0$ limit and the transverse integral in the $[k^{+}/p^{+}]_{+}$ term of Row V;
\item $R^{i}$ in Row II: $-2k^{i}$ should presumably be $-2q^{i}$;
\item the normalisation $\kappa$ of $[\rho,\rho]$, needed to combine the commutator terms with the existing two-$\rho$ terms.
\end{enumerate}

\section*{Numerical checks}
\begin{itemize}
\item Cancellation I: the colour identities were checked for $SU(2)$ with non-commuting, matrix-valued $\rho$ ($10^{-13}$).
\item Cancellation II: the three $\xi$-integrals (for $\mathcal{I}_{1}$, $\mathcal{I}_{2}$ and the finite part of $\mathcal{I}_{3}$) were checked by quadrature against $t_{16},\dots,t_{22}$ ($10^{-12}$ or better); the change of variables was checked pointwise; the split \eqref{eq:CVsplit} and the identities ``(i) $=-2\mathcal{W}_{\rm III}+\mathcal{Z}$'', ``(ii) $=+2\mathcal{W}_{\rm III}-\mathcal{Z}$'' were checked with matrix-valued $\rho$ ($10^{-14}$).
\item $S_{2}+S_{3}$ of the file equals $\sum_{16}^{21}t_{n}$ ($10^{-14}$).
\end{itemize}

\end{document}
